% Carried over verbatim from monograph/main.tex (Study 1 record, 13 Sep 2026).
\chapter{Conclusion}\label{s1:ch:ch18_conclusion}

This project asked whether a file-level train/test boundary corresponds to an independent acquisition boundary in public cervical cytology. It recovered and validated the hidden observation-to-parent hierarchy across four datasets, quantified the analytic and observed structural sharing, demonstrated literal source-pixel sharing in SIPaKMeD, and showed with a shuffled-label-controlled 424-way probe that parent identity remains strongly learnable after multiple cue-removal transformations. It then re-ran six classifiers and seven segmentation models over five fixed seeds and four matched protocols, preserving effect sizes, bootstrap intervals, multiplicity families, failures, and runtime metadata.

The strongest result is methodological: parent-disjoint evaluation is necessary whenever multiple released observations descend from one acquisition unit. The downstream neural results point in a consistent direction but do not justify a universal multiplicity-corrected accuracy-inflation claim. Four of six classifier architectures also violate the preregistered monotonic score ordering, a useful negative result showing that stricter grouping and lower finite-sample score are not synonymous.

CytoFair turns that lesson into reusable artifacts: provenance maps, duplicate components, fixed splits, hashes, fail-closed validators, run ledgers, and a raw-download structural pipeline. The next scientific step is a separately preregistered specimen/episode-disjoint fish-imagery V2; the next release step is a virgin-host clean-clone attestation. Until the latter, the numerical study is complete but the strongest release gate remains honestly pending.
